\section{Introduction}
\label{sec:introduction}

A judge orders the government to buy a cancer drug within
\BPurgentPlanLow--\BPurgentPlanHigh{} days. The procurement officer has
no time to aggregate demand across other purchases, no time to negotiate
with multiple suppliers, and faces personal fines if delivery fails.
The drug arrives. The natural question is why it costs more---and
whether the cost margin is best understood as suppliers extracting a
markup from a constrained buyer, or as a constrained buyer ending up
with a different, more expensive supplier altogether.

The answer matters for procurement design. If urgency raises prices
through within-firm markup, the policy lever is on contract terms---
deterring opportunistic pricing under sanction exposure. If urgency
forces the buyer onto a different point of the supplier distribution,
the lever is on demand aggregation and supplier search. The two readings
are observationally similar in standard data. They are observationally
distinct only in data, like the bid-level data we use here, that
allow comparison across regimes within the same firm-buyer-item triple.

This paper provides such evidence using all pharmaceutical procurement
by the S\~ao Paulo State Department of Health (SES/SP) from
\BPsampleStartYear{} through \BPsampleEndYear: \BPnPOIfull{}
purchase-offer-item observations conducted through the state's
electronic procurement platform (BEC). Court orders and administrative
requests partition urgent purchases into two types that share all
planning constraints---compressed timelines, small quantities, diverted
budget resources, the same BEC platform and bidding procedures---but
differ in penalty exposure. Litigated purchases expose officials to
fines, personal liability, and fund seizure if delivery fails.
Administrative requests, instituted by SES/SP in \BPsampleStartYear,
carry no such penalties. Comparing these two types within the same
item is what we call the ``under the gun'' (UTG) comparison.

The administrative subsample is itself selected by an SES/SP scientific
committee on cost-effectiveness criteria. Pre-period probit evidence
confirms the committee admits cases that are systematically less
expensive (probit coefficient on pre-period log reference price
$=-0.051$, $p<0.001$) and more clinical-essential (SUS-formulary
coefficient $-0.265$, $p<0.001$); rejected cases land in the litigated
sample. The naive UTG gap therefore confounds the sanction effect with
selection on item economics. We bound the selection wedge formally:
under Manski-Lee monotone restrictions on the committee rejection rule,
the UTG gap lies in $[\BPutgBoundLow, \BPutgBoundHigh]$. This range
remains economically meaningful; the rest of the paper asks where it
comes from.

Three findings frame the answer.

\emph{First}, the headline urgent-vs-ordinary premium is
\BPnegPctHeadline{} across the full sample, with the same item under
compressed timelines attracting \BPfirmsPctHeadline{} fewer bidding
firms. Urgent tenders nonetheless succeed \BPsuccessPP{} \emph{more}
often, ruling out an alternative in which urgency simply selects
easier purchases.

\emph{Second}, the price margin is \emph{not} a within-firm markup.
Within firm-buyer-item triples observed under both the litigated and
the administrative urgent regime (\BPutgTripleCountUTG{} triples,
\BPutgTripleNUTG{} observations), the Admin coefficient is
$\BPutgTripleCoefUTG$ (SE $\BPutgTripleSEUTG$, $p\!\approx\!0.4$). The
same firm sells the same item to the same buyer at essentially the
same per-unit price under either regime. Whatever pushes the
across-sample price up, it is not opportunistic pricing by suppliers
that already serve the buyer's procurement.

\emph{Third}, the bounded UTG admits an arithmetic decomposition into
a mechanical bulk-discount component (admin orders are $\sim\!3{.}3\times$
larger than litigated orders, mechanically pulling the admin-vs-lit
price gap to about $-0.397$ log-points), the within firm-buyer-item
null, and a positive supplier composition residual ($+0.103$
log-points). Same firms price the same items the same way regardless
of regime. Different firms win urgent contracts under the two regimes,
and the equilibrium-weighting shift is the residual we read as the
sourcing channel.

\paragraph{Three contributions.}
\emph{First}, we provide direct within-item evidence on the price
impact of court-mandated procurement, exploiting the institutional
within-urgent contrast unique to S\~ao Paulo. \emph{Second, and
conceptually most novel, we distinguish the sourcing channel of
procurement-cost variation from the pricing channel.} Standard
intuition under accountability mechanisms relies on within-firm pricing
distortion \citep{prendergast2007bureaucrats}; our within firm-buyer-item
triple evidence cleanly rejects that mechanism in this setting. The
empirical signature of accountability-induced inefficiency in our
data is mismatched suppliers and fragmented demand, not extorted
prices. \emph{Third}, in contrast to \citet{coviello2018court}---where
slow courts \emph{extend} procurement timelines and raise prices---we
study judicial injunctions that \emph{compress} timelines, with the
cost concentrated in items where demand aggregation matters most. The
two papers identify alternative institutional channels through which
judicial intervention raises procurement costs.

\paragraph{Welfare.}
Applying the bounded UTG to the \BPwelfareBoundAnnualSpend{} of
annual litigated spending in S\~ao Paulo, with an admissibility-share
calibration of \BPwelfareBoundAdmissibility, gives a per-unit-price
welfare bound of \BPwelfareBound{} per year, with the Lee-endpoint
range delivering $[\BPwelfareBoundCIlow, \BPwelfareBoundCIhigh]$. We
report a single bounded figure rather than the stacked ranges common
in the policy literature.

\paragraph{Position vis-\`a-vis the literature.}
The paper sits between the public-economics literature on procurement
frictions \citep{bandiera2009active, best2023procurement, bosio2022public,
decarolis2018buyer} and the institutional literature on the
judicialization of health in Latin America \citep{wang2015right,
ferraz2009right, biehl2009will}. It contributes to a smaller literature
on bureaucratic distortion under accountability mechanisms
\citep{rasul2018management, williams2021bureaucratic, bandiera2021allocation}
and on competition design under contractual incompleteness
\citep{carril2026competition}. Our methodological discipline avoids the
post-treatment-bias trap (\citealt{acharya2016explaining}): we read
quantity-controlled and firm-fixed-effect specifications as
channel-isolating, not as direct-effect estimators, and present the
within firm-buyer-item triple as the mechanism's positive evidence.

The remainder of the paper proceeds as follows.
Section~\ref{sec:institutional} describes the institutional setting.
Section~\ref{sec:data} presents the data. Section~\ref{sec:empirical}
details the empirical strategy and the selection-bound design.
Section~\ref{sec:results} presents results, organized around the
sourcing-vs-pricing pivot. Section~\ref{sec:conclusion} concludes.
